\chapter*{Introduzione}
\addcontentsline{toc}{chapter}{Introduzione}
\label{ch:introduction}

% Capitolo non numerato (fuori dalla sequenza Capitolo 1, 2, ...): le
% sezioni interne restano numerate ma senza prefisso di capitolo (1, 2, ...
% invece di 1.1, 1.2, ...), per non farle sembrare appartenere al Capitolo 1
% vero e proprio (Fondamenti matematici). Ripristinato in fondo al capitolo.
\let\oldthesection\thesection
\renewcommand{\thesection}{\arabic{section}}

\section{Il mercato delle criptovalute come sistema interconnesso}
\label{sec:crypto-market-interconnected}

Il mercato delle criptovalute appare, ad un primo sguardo, come un insieme ampio ed eterogeneo di asset indipendenti: migliaia di token, ciascuno con una propria tecnologia sottostante, una propria comunità di sviluppatori e un proprio caso d'uso dichiarato. Le serie storiche dei prezzi raccontano una storia diversa. I rendimenti delle criptovalute tendono a muoversi insieme, spesso con un'intensità superiore a quella osservabile tra azioni o materie prime di pari volatilità~\cite{corbet2019cryptocurrencies}. Se gli asset non sono indipendenti, un modello che li tratti come tali sta scartando per costruzione un'informazione potenzialmente rilevante.

Bitcoin resta il nodo dominante di questa rete di dipendenze: pur avendo perso quota di capitalizzazione rispetto ai primi anni del mercato, la sua traiettoria di prezzo continua ad anticipare quella della maggior parte degli \enquote{altcoin}\footnote{Con \enquote{altcoin} (da \emph{alternative coin}) si indica genericamente ogni criptovaluta diversa da Bitcoin.}, che mostrano tipicamente una sensibilità elevata ai suoi movimenti. Gli studi sulla connettività dinamica tra criptovalute, basati sul framework dello \emph{spillover index}\footnote{Indice proposto da Diebold e Yilmaz per misurare quanta parte della varianza dei rendimenti di un asset sia spiegata da shock — variazioni improvvise e inattese dei rendimenti — originati su altri asset, distinguendo chi genera shock netti (fonte) da chi li assorbe (ricevitore)~\cite{dieboldyilmaz2012spillover}.}, confermano questo ruolo: Bitcoin ed Ethereum emergono come le principali fonti nette di shock all'interno della rete di connessione tra asset, mentre la maggior parte degli altcoin si comporta da ricevitore netto~\cite{ji2019dynamic}.

L'evidenza più netta di questa interdipendenza emerge nei momenti di crisi, quando le correlazioni tra asset si intensificano invece di attenuarsi. Tre episodi recenti confermano questa dinamica. Nel maggio 2021 l'inasprimento delle restrizioni cinesi su scambi e mining provocò un crollo generalizzato del mercato, con perdite infragiornaliere superiori al 30\% sui principali asset~\cite{griffith2023cryptocurrency}. Un anno dopo, il collasso dell'ecosistema Terra/Luna, legato alla stablecoin algoritmica UST\footnote{Criptovaluta progettata per mantenere un valore stabile, tipicamente ancorato a una valuta come il dollaro statunitense, tramite riserve di collaterale o meccanismi algoritmici.}, cancellò gran parte della capitalizzazione dell'ecosistema in pochi giorni. Questo crollo ebbe conseguenze verificabili anche sulla struttura della rete di correlazione tra criptovalute ad alta capitalizzazione non direttamente coinvolte~\cite{briola2023anatomy}. Nel novembre 2022, il fallimento dell'exchange FTX produsse effetti analoghi sulla connettività tra asset privi di qualunque legame operativo diretto con l'exchange~\cite{akyildirim2023understanding}.

Questi episodi mostrano che le serie di prezzo delle criptovalute non sono indipendenti e che il grado di questa dipendenza non è stabile nel tempo. Le correlazioni tra asset si intensificano nei periodi di crisi e si allentano in quelli di relativa stabilità~\cite{onnela2003dynamics}, in linea con il \emph{volatility clustering} (\cref{sec:financial-time-series-returns})~\cite{cont2001empirical}. Trattare ogni criptovaluta come indipendente dalle altre comporta la perdita di un'informazione potenzialmente rilevante.

\section{Limiti dei modelli univariati}
\label{sec:limits-univariate-models}

La famiglia di modelli più diffusa per la previsione di serie storiche finanziarie affronta il problema in modo univariato: un modello AR viene addestrato separatamente su ciascun asset, utilizzando come unico input la storia passata dell'asset stesso~\cite{tsay2010analysis}. Questo approccio tratta ogni criptovaluta come un sistema chiuso, isolato dal resto del mercato: per costruzione, il modello non ha accesso ad alcuna informazione sulle altre serie e non può quindi rappresentare esplicitamente la dipendenza tra asset. Un modello univariato applicato a Bitcoin, ad esempio, non può in linea di principio \enquote{vedere} un segnale di stress che si sta propagando da un'altra parte del mercato, non essendo quel segnale tra i suoi input.

Questo limite si traduce direttamente nell'incapacità di anticipare i fenomeni di contagio discussi nella \cref{sec:crypto-market-interconnected}. Una collezione di modelli univariati indipendenti non può rappresentare l'intensificarsi delle correlazioni durante una crisi, perché tale intensificazione è una proprietà collettiva del sistema di asset, non un tratto osservabile nella singola serie presa da sola.

La statistica multivariata classica offre una risposta parziale a questo problema con il modello VAR, in cui il vettore dei rendimenti di tutti gli asset al tempo $t$ viene regredito sui propri valori ritardati~\cite{lutkepohl2005new}. A differenza dei modelli univariati, un VAR cattura esplicitamente le relazioni tra serie: ogni asset è messo in relazione con il passato di tutti gli altri, non solo con il proprio. Il prezzo di questa espressività è tuttavia alto. Per un sistema di $N$ asset e un ordine di ritardo $p$, il numero di parametri da stimare cresce come $O(N^2 p)$. Con poche decine di criptovalute il numero di parametri diventa rapidamente ingestibile rispetto alla quantità di dati storici disponibili. Ne derivano stime instabili e un rischio di overfitting severo, che rende il VAR classico poco praticabile su un universo ampio ed eterogeneo di asset. Esistono estensioni che attenuano il problema, come i VAR regolarizzati che impongono sparsità sui coefficienti~\cite{basu2015regularized}, ma restano ancorate ad una struttura di dipendenza globale e fissa, incapace di adattarsi ad una topologia di connessioni che cambia nel tempo.

I modelli univariati hanno quindi un costo che cresce linearmente con il numero di asset, ignorandone per costruzione la dipendenza reciproca; il VAR, invece, cattura questa dipendenza, ma con un costo che cresce quadraticamente. Occorre pertanto individuare una struttura capace di rappresentare le relazioni tra criptovalute ad un costo computazionale più sostenibile.

\section{Il grafo come intuizione strutturale}
\label{sec:why-graph-intuition}

Il mercato delle criptovalute, inteso come sistema di asset connessi, si formalizza come grafo pesato: ciascuna criptovaluta è rappresentata da un nodo, ogni coppia di asset è collegata da un arco il cui peso misura l'intensità della relazione tra le due serie, e a ciascun nodo è associato un insieme di feature osservate nel tempo, come rendimento e volume.

Il grafo completo, con un arco per ogni coppia di asset, è però denso e per questo poco trattabile. L'econofisico Mantegna propose per primo un metodo per isolare, a partire dalla matrice di correlazione, le relazioni più significative tra gli asset: un albero di copertura minimo (\emph{minimum spanning tree}, MST) che consente di ottenere una rappresentazione sintetica della struttura di dipendenza del mercato~\cite{mantegna1999hierarchical}. Tecniche successive, sviluppate da Tumminello, raffinano questo filtraggio conservando un numero maggiore di connessioni rilevanti rispetto al solo albero minimo, restituendo una rappresentazione più ricca della struttura di correlazione originale~\cite{tumminello2005tool}.

Mantenere solo le connessioni più rilevanti, principio comune a queste due tecniche, risolve la tensione tra espressività e costo quadratico. Un grafo può restare sparso, includendo un arco solo tra le coppie di asset la cui relazione supera una soglia prefissata, piuttosto che stimare un coefficiente per ogni possibile coppia come richiesto dal VAR completo. Il costo computazionale e statistico dipende quindi esclusivamente dal numero di connessioni effettivamente presenti. Ogni nodo, inoltre, pur conservando una propria dinamica individuale, scambia informazioni con i nodi a cui è connesso attraverso gli archi del grafo. Questa rappresentazione selettiva delle dipendenze, tuttavia, perde di efficacia nei momenti di crisi, quando le connessioni si moltiplicano e la distinzione tra relazioni rilevanti e trascurabili tende a svanire. Il grafo va quindi inteso come una struttura dinamica, con densità variabile nel tempo.

Rappresentare il mercato come un grafo non è però sufficiente a renderlo predittivo: descriverne la topologia in un dato istante non equivale ad un meccanismo per apprendere da questa struttura e trasformarla in una previsione. È necessaria un'architettura capace di elaborare il segnale di ogni nodo tenendo conto, in modo esplicito, dell'assetto del grafo che lo circonda.

\section{Domande di ricerca e struttura della tesi}
\label{sec:objectives-research-questions}

La tesi si articola attorno a due questioni aperte, emerse dalle sezioni precedenti: quale vantaggio predittivo offra la struttura relazionale delle criptovalute rispetto ad un approccio che le tratti come asset indipendenti e come si trasformi la topologia del grafo nel passaggio da un regime di mercato stabile a uno di crisi.

Dal punto di vista teorico, la tesi affronta la prima questione costruendo l'apparato matematico necessario a definire una Graph Convolutional Network in grado di sfruttare la struttura relazionale discussa nella \cref{sec:why-graph-intuition}. Sul piano empirico, la prima questione trova risposta nel \cref{ch:from-cnn-to-gnn}, in cui le prestazioni della GCN sono confrontate con quelle delle baseline di riferimento, e nel \cref{ch:backtest}, che ne valuta il vantaggio sul piano economico. La seconda questione viene trattata nel \cref{ch:topological-evolution}, che analizza l'evoluzione della struttura relazionale attorno a tre eventi di crisi.

La tesi è organizzata come segue. Una breve panoramica precede il primo capitolo, presentando l'organizzazione del progetto software su cui si basano i riferimenti al codice nei capitoli successivi. Il \Cref{ch:graph-laplacian-theory} introduce gli strumenti di teoria dei grafi necessari allo studio. Il \Cref{ch:from-cnn-to-gnn} deriva l'architettura della Graph Convolutional Network di Kipf \& Welling e le baseline con cui viene confrontata, coprendo anche il protocollo di validazione walk-forward con cui vengono valutate e le insidie metodologiche note in letteratura che quel protocollo previene. Il \Cref{ch:financial-series-and-graph-construction} formalizza la costruzione del grafo dinamico di correlazione a partire dalle serie storiche dei rendimenti, dalla scelta dell'universo di asset fino al sostrato usato dalla GCN. Il \Cref{ch:topological-evolution} applica gli strumenti spettrali alla serie storica del grafo così costruito, misurandone l'evoluzione attorno a tre eventi di crisi reali. Il \Cref{ch:backtest} valuta le previsioni sul piano economico, con una strategia di trading elementare al netto dei costi di transazione. Il \Cref{ch:conclusions} sintetizza infine i risultati rispetto alle due questioni presentate, discutendone limiti e possibili sviluppi futuri.

\let\thesection\oldthesection
